% ============================================================
% SECTION 8 - CONCLUSION
% ============================================================
\section{Conclusion}
\label{sec:conclusion}

%------------------------------------------------------------
\subsection{Summary}
\label{sec:conclusion:summary}
%------------------------------------------------------------

We identified a structural limitation of the pooled EnbPI algorithm
\citep{xu2023conformal} when it is applied to seasonally
heteroskedastic time series: pooled calibration introduces a
persistent, season-dependent coverage bias of magnitude $2b_{s^*}$
that does not vanish as the training sample grows
(Proposition~\ref{prop:pooled_gap}).  Two schemes remove it, and they
are the endpoints of the one-parameter family~\eqref{eq:family}:
stratifying the calibration buffer by season, which is
\texttt{EnbPI-S}, and standardising each residual by a season-specific
scale before pooling, which is \texttt{EnbPI-N}.  The second is the
more efficient of the two, delivering a comparable coverage level from
an interval about a fifth narrower, and it is valid only under the
scale-mixture restriction of Assumption~\ref{ass:seasonal}; the first
does not need that restriction, and when the seasons differ in the
shape of the error distribution rather than in its scale it is the only
one of the two that equalises coverage, cutting the seasonal spread to
a third of what pooling leaves while normalisation leaves it
essentially untouched.  A permutation test on the training residuals
decides between the two configurations, and on both Brazilian series it
rejects the pure-scale restriction.

The paper's second finding was not the one it set out to make.  The
coverage that stratified calibration loses at the per-season sample
sizes a monthly series provides, eleven percentage points at
$\Ts = 20$, is not the price of conditioning on the season.  It is the
price of two properties of the calibration step of the original
algorithm: the empirical quantile of a short buffer is biased inwards,
and the minimum-width line search selects the narrowest of many
candidate quantile pairs computed from the same residuals, which alone
costs four to six percentage points at these buffer lengths.  Building
the interval from conformal order statistics with the asymmetry fixed
at $\alpha/2$ repairs both, and stratified coverage at twenty residuals
per season then attains its nominal level in all four seasonal designs
we study.  What remains is a cost in width rather than in coverage, and
it is the honest price of a finite-sample guarantee on a short buffer;
blending towards the normalised endpoint recovers about a tenth of it.

On the theoretical side, we proved that the per-season coverage gap of
\texttt{EnbPI-S} decays at rate $O(\sqrt{\log T_s / T_s})$, where
$T_s = T/\Sper$ is the per-stratum training size, and we derived an
asymptotic dominance result (Proposition~\ref{prop:strat_dominance})
together with a finite-sample condition~\eqref{eq:dominance_condition}
that links the benefit of stratification to the size of $b_{s^*}$ and
to the training sample length, although that condition turns out to be
too conservative to be checked at any sample size an economic series
provides.  Proposition~\ref{prop:family_coverage} bounds the coverage
of the interpolated interval and shows that asymptotic validity
requires the weight to tend to one whenever the shape discrepancy is
nonzero.  The result that carries the most weight, however, is the
smallest: Proposition~\ref{prop:finite_sample} records that the
conformal endpoints deliver at least $1-\alpha$ exactly, at every
buffer length above $2/\alpha - 1$, and no more than
$1-\alpha + 2/(\Ts+1)$.  That pair of bounds is what the repaired
construction is doing, and it explains both halves of what the
experiments show: the nominal coverage at twenty residuals per season
and the conservatism that comes with it.  Two limits of that theory should be read alongside it.  It is
stated for a fixed quantile pair, whereas the algorithm chooses the
pair from the same buffer, and Experiment E6 measures what that costs;
and it says nothing about the data-driven weight, for the reason given
in Remark~\ref{rem:lambda_honesty}.

Empirically, the two applications share the test window January 2015 to
December 2024.  On food-at-home IPCA ($\Ts = 20$), the stratified
scheme built from conformal order statistics attains exactly $90.0\%$
coverage with no dispersion across the twelve calendar months, against
$79.2\%$ for the same buffer read as an empirical quantile.  On
Brazilian export growth ($\Ts = 35$, BCB series 1402), the same change
raises the worst month from $70\%$ to $90\%$ and overall coverage from
$86.7\%$ to $95.0\%$, the overshoot being the conservatism of an
interval whose endpoints are still the extremes of a
thirty-five-observation buffer, at roughly twice the width; blending
towards the normalised endpoint brings the width back by a fifth at the
same coverage.  In both applications the seasonal width pattern is
economically interpretable, widest in June for food prices and in April
for exports, at the peak of the soybean harvest-export season.

%------------------------------------------------------------
\subsection{Limitations}
\label{sec:conclusion:limitations}
%------------------------------------------------------------

The main practical limitation is not the one earlier drafts of this
paper reported.  It is not that a season buffer of twenty residuals
cannot support a valid interval, since Experiment E6 shows that it can;
it is that the interval it supports is wide.  Between the two
thresholds discussed after~\eqref{eq:feasible} the conformal endpoints are the
extremes of the buffer, so at $\Ts = 20$ the stratified interval is
about forty per cent wider than the pooled one, and the excess is only
partly recovered by blending towards the normalised endpoint.  An
analyst who needs a narrow interval more than a season-conditional
guarantee should still pool, and should know that is the trade being
made.

The requirement interacts directly with data frequency and
structural stability.  For \emph{monthly} data ($\Sper = 12$), the
threshold $T_s \gtrsim 30$ translates to $T \gtrsim 360$ months,
30 years, which is achievable for standard macroeconomic series such
as industrial production or price indices.  For \emph{quarterly} data
($\Sper = 4$), the same $T_s$ threshold requires only
$T \gtrsim 120$ quarters, again 30 years, though the assumption of
structural stability over three decades is considerably harder to
justify for quarterly aggregates such as GDP: the 2008--09 financial
crisis and the 2020 COVID recession both constitute substantial
regime shifts that can invalidate a calibration buffer spanning that
horizon.  The practical implication is that \texttt{EnbPI-S} is most
naturally suited to \emph{monthly or higher-frequency} data, where a
large enough $T_s$ is reached within a single economic cycle, where
seasonal het\-ero\-sked\-as\-tic\-ity is typically more pronounced
(harvest and weather effects operate at the monthly scale), and where
structural breaks are less likely to contaminate the entire
calibration window.  For purely quarterly series, practitioners may
prefer to reduce $\Sper$ by grouping adjacent quarters into broader
``seasons'', or to use a shorter rolling-window calibration.

The framework also assumes a known, fixed seasonal period $\Sper$.
In practice the relevant cycle may be quarterly ($\Sper = 4$),
weekly ($\Sper = 52$), or irregular, as with fiscal years that do
not align with calendar years.  The algorithm extends immediately to
any integer $\Sper$, but the selection of $\Sper$ is left to the
practitioner.

Structural breaks pose a further limitation.  The sliding-window
update adapts to slow drift in the seasonal volatility structure,
though it is not designed to handle abrupt breaks.  Part of the
coverage shortfall in the food-IPCA application plausibly reflects
such a break, since COVID-era inflation dynamics (2020--22) lie
partly outside the support of the 1995--2014 training distribution.
Robust recalibration after detected breaks remains an open problem.

%------------------------------------------------------------
\subsection{Future Directions}
\label{sec:conclusion:future}
%------------------------------------------------------------

\citet{xu2023sequential} proposed SPCI, which exploits residual
autocorrelation by training a quantile random forest on the history
of residuals.  \texttt{EnbPI-S} and SPCI address orthogonal sources
of miscoverage, since SPCI corrects for serial dependence and
\texttt{EnbPI-S} corrects for seasonal scale variation.  A combined
method could therefore apply the quantile-forest step within each
seasonal buffer, potentially achieving intervals that are at once
autocorrelation-adjusted and season-calibrated.  We leave the
formalisation of such a combined estimator to future work.

Another direction is adaptive stratum selection.  Rather than
stratifying by calendar month, one could stratify by a data-driven
cluster of months with similar residual distributions; a Gaussian
mixture model or a $k$-means clustering of the monthly empirical
quantile vectors would allow the number of strata to be chosen from
data, trading off stratification bias against per-stratum sample
size.

A third direction concerns the interior of the family~\eqref{eq:family}.
Experiment E6 shows that blending the stratified endpoint towards the
normalised one recovers about a tenth of the width that a finite-sample
guarantee costs on a short buffer, at essentially unchanged coverage,
but also that choosing the weight from data does not beat fixing it at
one half.  A weight that is worth estimating would have to exploit
something the cross-validated rule does not, most plausibly the fact
that the seasons are not twelve unrelated problems: a hierarchical
estimator that shrinks each season towards a common seasonal profile
rather than towards the pooled buffer is the natural next step, and it
is also the setting in which shrinkage across many related quantities
has a theory to lean on.  Whether such an estimator can be given a
finite-sample guarantee, given that the weight would be a function of
the same residuals that form the buffer, is the harder half of that
question.

Finally, the framework applies directly to quarterly, daily or
intraday seasonality by adjusting $\Sper$.  For high-frequency data
with large $\Sper$, as in hourly electricity demand with an annual
cycle ($\Sper = 8760$), the per-stratum sample sizes would be very
small and alternative grouping strategies would become necessary; the
theoretical analysis extends to this regime with minor adjustments to
the DKW-based argument in Appendix~\ref{app:proofs}.

Beyond the specific algorithm, the broader lesson of the exercise is
that, for seasonal economic series, the calibration step deserves at
least as much attention as the point-forecast model: the
SARIMA benchmark of Section~\ref{sec:simulations} loses more than
twenty percentage points of coverage in its worst month even though
its conditional mean closely approximates the true one.
